\subsection{Fungsi Invers}
\begin{definition}[Fungsi Invers]
  Fungsi invers adalah fungsi yang ``Membalik'' fungsi asalnya.
  \begin{itemize}
  \item Jika fungsi asal memetakan $x \implies y$.
  \item Maka fungsi invers memetakan $y \implies q$.
  \end{itemize}
  Notasi fungsi invers adalah:
  \[ f^{-1}(x) \]
  Bukan pangkat negatif tapi lambang kebalikan. Dibaca ``$f$ invers dari $x$''.
\end{definition}

\subsubsection{Menetukan Fungsi Invers}
\begin{enumerate}
\item Misalkan $y = f(x)$
\item Tukar Posisi $x$ dan $y$. ($x = f(y)$).
\item Ubah kembali ke $y$ = $f^{-1}(x)$.
\end{enumerate}

\begin{example}[Fungsi Invers]
  \hfill \\
  Tentukan fungsi invers dari fungsi
  \[ f(x) = 2x + 3 \]
  $f(x) = 2x + 3$ \\
  $y = 2x + 3$ \\
  $-2x = -y + 3$ \\
  $2x = y - 3$ \\
  $x = \frac{y - 3}{2}$ \\
  $f(y) = \frac{y - 3}{2}$ \\
  $f^{-1}(x) = \frac{x - 3}{2}$
\end{example}
